\section{Conclusiones}

El análisis del fenotipo \textit{Abnormal drinking behaviour} (HP:0030082) mediante un enfoque de biología de sistemas ha permitido identificar patrones estructurales y funcionales que explican cómo la regulación de la ingesta de líquidos emerge de la interacción entre múltiples vías moleculares. Gracias al uso de la \textit{Human Phenotype Ontology} (HPO) como punto de partida estandarizado, fue posible construir una red génica coherente y explorar su organización interna a nivel topológico y fisiológico.

\subsection{Principales hallazgos del análisis de red}
El estudio de las interacciones proteicas reveló una arquitectura molecular organizada en torno a genes clave relacionados principalmente con la función renal, la señalización hormonal y la neurotransmisión. La red final, compuesta por 38 nodos y 112 aristas, mostró una \textbf{distribución de grados heterogénea} y una modularidad consistente con lo esperado en redes biológicas complejas.

Entre los hallazgos más relevantes destacan:
\begin{itemize}
    \item \textbf{Predominio etiológico renal}: Genes como AQP2, BSND, PKHD1, SLC12A3 y NPHP1 presentan posiciones estructurales centrales, reflejando que la capacidad renal para modular el transporte de agua y electrólitos constituye un componente clave del fenotipo analizado.

    \item \textbf{Hubs críticos con relevancia fisiológica}: La centralidad de genes como AVPR2, ACE, OPRM1, DRD2 y GABRA2 demuestra que el fenotipo no depende únicamente del riñón, sino que integra mecanismos hormonales, conductuales y neuroendocrinos.
    
    \item \textbf{Estructura funcional y modularidad}: La presencia de clústeres bien definidos confirma que los genes asociados al fenotipo se organizan en módulos coherentes, vinculados con procesos como la homeostasis hídrica, la señalización hormonal y el control del comportamiento ingestivo.
\end{itemize}

En conjunto, estos resultados refuerzan la idea de que los patrones de ingesta anormal no derivan de una única causa molecular, sino de la interacción de varios sistemas interdependientes.

\subsection{Validación de la hipótesis y del marco conceptual inicial}
La hipótesis planteaba que los genes asociados al fenotipo constituían un módulo funcionalmente coherente. Los resultados obtenidos permiten afirmar que: 
\begin{itemize} 
    \item \textbf{1. }La centralidad de AVPR2, AQP2 y otros genes clave confirma la integración funcional entre vías osmóticas, hormonales y conductuales. 
    \item \textbf{2. }La modularidad de la red, con clústeres bien definidos, coincide con el concepto de “módulos biológicos” descrito por Barabási. 
    \item \textbf{3. }El enriquecimiento funcional GO:BP refuerza la relevancia fisiológica del conjunto. 
    \item \textbf{4. }La topología de la red revela un comportamiento típico de redes biológicas complejas, no de agrupaciones arbitrarias.
\end{itemize} 

Por tanto, se valida el marco conceptual expuesto en la Introducción: \textbf{el fenotipo \textit{Abnormal drinking behaviour} no emerge de un gen aislado, sino de la desregulación integrada de varios sistemas moleculares interdependientes}.

\subsection{Limitaciones y perspectivas futuras}
Pese a la solidez de los hallazgos, deben considerarse varias limitaciones:
\begin{itemize}
    \item El tamaño relativamente reducido del conjunto génico limita la potencia estadística, especialmente en análisis de enriquecimiento.
    \item STRING, aunque exhaustivo, no cubre la totalidad de las interacciones biológicas posibles.
    \item La red representa asociaciones estáticas; no incorpora dinámica temporal ni especificidad tisular.
    \item Algunos procesos altamente especializados (cilios renales, regulación fina de la sed hipotalámica) pueden quedar infrarrepresentados en ontologías generales como GO.
\end{itemize}

Para profundizar en la caracterización del fenotipo, futuros trabajos podrían integrar:
\begin{itemize}
    \item Datos transcriptómicos o proteómicos específicos de tejidos relevantes (riñón, hipotálamo, neuronas dopaminérgicas).
    \item Modelos dinámicos y simulaciones de señalización.
    \item Ánálisis de variantes genéticas patológicas.
    \item Redes multimodales que incorporen factores ambientales, conductuales y hormonales.
\end{itemize}